% GENERATED FILE. Edit canonical lessons, then run npm run generate:books.
% canonical-source-hash: fnv1a64:24dc80513ac69bcf
% canonical-lessons: TE-C505-jali, TE-C505-karuna, TE-C505-adarincu, TE-C505-aduko, TE-C505-samudayincu

\chapter{Kindness and Comfort}
\label{ch:te-kindness-and-comfort}
\hlchaptermodality{505}

\begin{chapteropening}
\textbf{By the end of this chapter:} \emph{I can say that someone showed pity or kindness and gave comfort.}

Complete the last lesson of chapter 505: I can say that someone showed pity or kindness and gave comfort.
\end{chapteropening}

\section[\texorpdfstring{\te{జాలి} (jāli)}{జాలి (jāli)}]{\te{జాలి} (jāli) — pity; sympathy}

\label{lesson:TE-C505-jali}



\begin{quote}
Before the new one: say the Telugu for to appreciate, then the Telugu for acclaim.
\end{quote}

\subsection*{You'll want to know: \te{జాలి}}
\textbf{\te{జాలి}} — \emph{jāli} — \textquotedblleft{}pity; sympathy\textquotedblright{}.

\begin{quote}
\te{ఆకలితో} \te{ఉన్న} \te{కుక్కపిల్లను} \te{చూసి} \te{నాకు} \te{జాలి} \te{వేసింది}.
\end{quote}

\begin{quote}
\emph{ākalitō unna kukkapillanu cūsi nāku jāli vēsindi.}
\end{quote}

\begin{quote}
Seeing the hungry puppy, I felt pity.
\end{quote}

\te{జాలి} \te{వేయు} or \te{జాలి} \te{కలుగు}, to feel pity; -\te{మీద} \te{జాలిపడు}, to pity someone.

\begin{cousinweb}
\textbf{\te{జాలి}} (\emph{jāli}) is a Telugu word for the pity you feel for someone weaker; its origin is not settled.
\end{cousinweb}

\begin{grammarlens}[title={What you've built}]
One more word for describing what you lived and felt.
\end{grammarlens}

\subsection*{Your turn}
\begin{itemize}
\raggedright
  \item \emph{Say it:} \emph{jāli}
  \item \emph{Say it:} \emph{jāli}, once more
  \item \emph{Say it:} \emph{ākalitō unna kukkapillanu cūsi nāku jāli vēsindi.} — Seeing the hungry puppy, I felt pity.
  \item \emph{Recall:} say the Telugu for to appreciate, then the Telugu for acclaim, then say \emph{jāli} again
\end{itemize}

\begin{tcolorbox}[breakable,colback=teal!4,colframe=teal!35!black,title={Before you move on}]
What does \te{జాలి} mean? (\textquotedblleft{}Pity; sympathy\textquotedblright{}.) Say it once more.
\end{tcolorbox}
\section[\texorpdfstring{\te{కరుణ} (karuṇa)}{కరుణ (karuṇa)}]{\te{కరుణ} (karuṇa) — compassion; mercy}

\label{lesson:TE-C505-karuna}



\begin{quote}
Before the new one: say the Telugu for acclaim, then the Telugu for pity.
\end{quote}

\subsection*{You'll want to know: \te{కరుణ}}
\textbf{\te{కరుణ}} — \emph{karuṇa} — \textquotedblleft{}compassion; mercy\textquotedblright{}.

\begin{quote}
\te{ఆ} \te{డాక్టరు} \te{పేదల} \te{పట్ల} \te{కరుణ} \te{చూపించారు}.
\end{quote}

\begin{quote}
\emph{ā ḍākṭaru pēdala paṭla karuṇa cūpiñcāru.}
\end{quote}

\begin{quote}
That doctor showed compassion to the poor.
\end{quote}

\te{కరుణ} \te{చూపించు}, to show mercy; a formal word.

\begin{cousinweb}
\textbf{\te{కరుణ}} (\emph{karuṇa}) is borrowed from Sanskrit \emph{karuṇā}, \textquotedblleft{}pity\textquotedblright{} (CDIAL 2811).
\end{cousinweb}

\begin{grammarlens}[title={What you've built}]
One more word for describing what you lived and felt.
\end{grammarlens}

\subsection*{Your turn}
\begin{itemize}
\raggedright
  \item \emph{Say it:} \emph{karuṇa}
  \item \emph{Say it:} \emph{karuṇa}, once more
  \item \emph{Say it:} \emph{ā ḍākṭaru pēdala paṭla karuṇa cūpiñcāru.} — That doctor showed compassion to the poor.
  \item \emph{Recall:} say the Telugu for acclaim, then the Telugu for pity, then say \emph{karuṇa} again
\end{itemize}

\begin{tcolorbox}[breakable,colback=teal!4,colframe=teal!35!black,title={Before you move on}]
What does \te{కరుణ} mean? (\textquotedblleft{}Compassion; mercy\textquotedblright{}.) Say it once more.
\end{tcolorbox}
\section[\texorpdfstring{\te{ఆదరించు} (ādariñcu)}{ఆదరించు (ādariñcu)}]{\te{ఆదరించు} (ādariñcu) — to welcome warmly; to cherish}

\label{lesson:TE-C505-adarincu}



\begin{quote}
Before the new one: say the Telugu for pity, then the Telugu for compassion.
\end{quote}

\subsection*{You'll want to know: \te{ఆదరించు}}
\textbf{\te{ఆదరించు}} — \emph{ādariñcu} — \textquotedblleft{}to welcome warmly; to cherish\textquotedblright{}.

\begin{quote}
\te{కొత్తగా} \te{వచ్చిన} \te{కోడలిని} \te{అత్త} \te{ఆదరించింది}.
\end{quote}

\begin{quote}
\emph{kottagā vaccina kōḍalini atta ādariñcindi.}
\end{quote}

\begin{quote}
The mother-in-law warmly welcomed her new daughter-in-law.
\end{quote}

Past (I …): \te{ఆదరించాను}, \emph{ādariñcānu}.

The noun \te{ఆదరణ} is a warm reception or support.

\begin{cousinweb}
\textbf{\te{ఆదరించు}} (\emph{ādariñcu}) is built on Sanskrit \emph{ādara}, \textquotedblleft{}respect\textquotedblright{} (CDIAL 1142), with the Telugu verb ending -\te{ఇంచు}.
\end{cousinweb}

\begin{grammarlens}[title={What you've built}]
One more word for describing what you lived and felt.
\end{grammarlens}

\subsection*{Your turn}
\begin{itemize}
\raggedright
  \item \emph{Say it:} \emph{ādariñcu}
  \item \emph{Say it:} \emph{ādariñcu}, once more
  \item \emph{Say it:} \emph{kottagā vaccina kōḍalini atta ādariñcindi.} — The mother-in-law warmly welcomed her new daughter-in-law.
  \item \emph{Recall:} say the Telugu for pity, then the Telugu for compassion, then say \emph{ādariñcu} again
\end{itemize}

\begin{tcolorbox}[breakable,colback=teal!4,colframe=teal!35!black,title={Before you move on}]
What does \te{ఆదరించు} mean? (\textquotedblleft{}To welcome warmly; to cherish\textquotedblright{}.) Say it once more.
\end{tcolorbox}
\section[\texorpdfstring{\te{ఆదుకో} (ādukō)}{ఆదుకో (ādukō)}]{\te{ఆదుకో} (ādukō) — to stand by; to help in need}

\label{lesson:TE-C505-aduko}



\begin{quote}
Before the new one: say the Telugu for compassion, then the Telugu for to welcome warmly.
\end{quote}

\subsection*{You'll want to know: \te{ఆదుకో}}
\textbf{\te{ఆదుకో}} — \emph{ādukō} — \textquotedblleft{}to stand by; to help in need\textquotedblright{}.

\begin{quote}
\te{కష్టకాలంలో} \te{స్నేహితులు} \te{మమ్మల్ని} \te{ఆదుకున్నారు}.
\end{quote}

\begin{quote}
\emph{kaṣṭakālaṁlō snēhitulu mammalni ādukunnāru.}
\end{quote}

\begin{quote}
In hard times our friends stood by us.
\end{quote}

Past (I …): \te{ఆదుకున్నాను}, \emph{ādukunnānu}.

Help given when someone is in trouble, not everyday help.

\begin{cousinweb}
\textbf{\te{ఆదుకో}} (\emph{ādukō}) is a Telugu verb ending in the reflexive -\te{కో}; the origin of its stem is not settled.
\end{cousinweb}

\begin{grammarlens}[title={What you've built}]
One more word for describing what you lived and felt.
\end{grammarlens}

\subsection*{Your turn}
\begin{itemize}
\raggedright
  \item \emph{Say it:} \emph{ādukō}
  \item \emph{Say it:} \emph{ādukō}, once more
  \item \emph{Say it:} \emph{kaṣṭakālaṁlō snēhitulu mammalni ādukunnāru.} — In hard times our friends stood by us.
  \item \emph{Recall:} say the Telugu for compassion, then the Telugu for to welcome warmly, then say \emph{ādukō} again
\end{itemize}

\begin{tcolorbox}[breakable,colback=teal!4,colframe=teal!35!black,title={Before you move on}]
What does \te{ఆదుకో} mean? (\textquotedblleft{}To stand by; to help in need\textquotedblright{}.) Say it once more.
\end{tcolorbox}
\section[\texorpdfstring{\te{సముదాయించు} (samudāyiñcu)}{సముదాయించు (samudāyiñcu)}]{\te{సముదాయించు} (samudāyiñcu) — to console; to pacify}

\label{lesson:TE-C505-samudayincu}



\begin{quote}
Before the new one: say the Telugu for to welcome warmly, then the Telugu for to stand by.
\end{quote}

\subsection*{You'll want to know: \te{సముదాయించు}}
\textbf{\te{సముదాయించు}} — \emph{samudāyiñcu} — \textquotedblleft{}to console; to pacify\textquotedblright{}.

\begin{quote}
\te{పరీక్ష} \te{తప్పిన} \te{తమ్ముడిని} \te{అక్క} \te{సముదాయించింది}.
\end{quote}

\begin{quote}
\emph{parīkṣa tappina tammuḍini akka samudāyiñcindi.}
\end{quote}

\begin{quote}
Big sister consoled our little brother, who had failed the exam.
\end{quote}

Past (I …): \te{సముదాయించాను}, \emph{samudāyiñcānu}.

Calming someone who is upset, often a child.

\begin{cousinweb}
\textbf{\te{సముదాయించు}} (\emph{samudāyiñcu}) is a Telugu verb for calming and consoling someone; its origin is not settled.
\end{cousinweb}

\begin{grammarlens}[title={What you've built}]
One more word for describing what you lived and felt.
\end{grammarlens}

\subsection*{Your turn}
\begin{itemize}
\raggedright
  \item \emph{Say it:} \emph{samudāyiñcu}
  \item \emph{Say it:} \emph{samudāyiñcu}, once more
  \item \emph{Say it:} \emph{parīkṣa tappina tammuḍini akka samudāyiñcindi.} — Big sister consoled our little brother, who had failed the exam.
  \item \emph{Recall:} say the Telugu for to welcome warmly, then the Telugu for to stand by, then say \emph{samudāyiñcu} again
\end{itemize}

\begin{tcolorbox}[breakable,colback=teal!4,colframe=teal!35!black,title={Before you move on}]
What does \te{సముదాయించు} mean? (\textquotedblleft{}To console; to pacify\textquotedblright{}.) Say it once more.
\end{tcolorbox}
